\section{Contribuições}

Esta pesquisa contribui para a compreensão da interatividade em ambientes de sala de aula virtual, tendo como parâmetro o contexto da metodologia expositiva, e servindo como norteadora de tendências para esse cenário. Por meio do ESPEON, essa avaliação torna-se acessível e prática, em razão da baixa barreira de entrada para uso e implementação da ferramenta. Além disso, o monitoramento ocorre de forma transparente, em \textit{background}, minimizando desconfortos ao usuário durante o processo. 

Os relatórios de aula desempenham papel central na pesquisa, por sua flexibilidade de formato — podendo se apresentar tanto como um documento PDF simples quanto como relatórios mais elaborados, gerados a partir do consumo da \textit{API} de relatórios do sistema. Essa versatilidade facilita a condução de estudos futuros e a consolidação de dados provenientes de diferentes contextos, possibilitando análises mais abrangentes e precisas sobre o fenômeno investigado.
